\scanpage{196}
\fbox{\shortstack{Continuité de la\\Torsion}}

\diagram{\begin{tikzcd}[ampersand replacement=\&,row sep=small]G\arrow[d,dash]\\ S\end{tikzcd}}
schéma en groupes \inserted{comm.} projectif sur $S$ \deleted{\unsure{loc.|loc}} noethérien \uline{local}
\deleted{$G^0$ univ. ouvert sur $S$.}
$s\in S$ \inserted{pt fermé}, $s_i$ générisations maximales.
Pour tt $t\in S$, soit $G_t^0$ muni de la structure \uline{abélienne} induite, soit $T_t(G)=G_t/G_t^0$
et $\sigma(t)=\mathrm{rang}_{k(t)}T_t(G)$

Implications

(i) \deleted{$G$ \uline{plat} sur $S$,} $G^{\deleted{\unsure{0|\circ}}}$ \deleted{contient}\inserted{admet} un sous-schéma \uline{abélien} $A$ ayant
$G^0$ comme sous-ens. sous-jacent, et $G$ plat sur $S$

\margin{$\Updownarrow$}(ii) $G$ admet un $A$ comme dessus, et $T(G)\simeq G/A$ est
\uline{plat} sur $S$ [N.B. il est déjà fini sur $S$]

\margin{$\Downarrow$}(iii) Pour tout $t\in S$, on a \deleted{rang} $\sigma(s)\deleted{\unsure{\geq}}=\sigma(\deleted{\unsure{s}}t)$ \deleted{[\unsure{ou seul.} $\sigma(s)\leq\sigma(s_i)$]}

\margin{$\Updownarrow$}(iv) Pour tt \deleted{\unsure{$t\in S$, on}} \unsure{$i$, on a} $\sigma(\unsure{s|t})=\sigma(s_i)$ [ou seulemt $\sigma(s)\leq\sigma(s_i)$]

Si $G^0$ univ. ouvert sur $S$, \inserted{i.e. les fibres de $G$ ont ttes $=$ dim.} on a \textup{(iii)}$\Longleftrightarrow$\textup{(iv)}.

Si \inserted{de plus} $S$ est \uline{réduit}, \uuline{ttes} les conditions envisagées sont équivalentes.

On sait que (\unsure{théorie|Théorie} des \unsure{passages|groupes} au quotient \inserted{plat}) que si $A$
existe, alors $G\longrightarrow G/A$ est fid. plat et projectif, et que
alors $G$ plat sur $S$ sss $G/A$ plat sur $S$, ce qui donne
l'équivalence de (i) et (ii). De plus, si $A$ existe, alors
on a pour tt $t\in S$, $T_t(G)\simeq T(G)_t$ où $T(G)=G/A$.
Donc (ii) $\Longrightarrow$ (iv), \unsure{à|:|a} fortiori (ii) $\Longrightarrow$ (iii).

Pour (iii) $\Longleftrightarrow$ (iv), on \deleted{\unsure{aura besoin du}}\inserted{applique le}

\uline{Lemme} \inserted{Supposons $G^0$ univ. ouvert sur $S$, i.e. les fibres de $G$ ont ttes $=$ dim.} La fonction $\deleted{\unsure{f}}\sigma$ est \uline{semi-continue} \uline{supérieurement}.
Comme elle est évidemment constructible (cf l'argument
précédent \ldots\,.) on est ramené à prouver, dans le cas où $S$ est
le spectre d'un anneau de val. discrète, $s'$ son pt gén.,
que $\sigma(s)\geq\sigma(s')$. Soit $G^0$ muni de la structure adhérence
de celle de $G^0_{s'}$, donc $G^0$ est un \uline{ss-groupe} de $G$ \uline{plat}$/S$,
et on peut former le groupe quotient $Q=G/G_0$, qui est
fini sur $S$. Alors
\[
\sigma(s)\geq\mathrm{rang}\,[Q_s:k(s)]\geq\mathrm{rang}\,[Q_{s'}:k(\unsure{s|s'})]=\sigma(s')
\]
d'où l'inégalité voulue. Noter que la première inégalité est
une égalité sss $G_s^0$ est réduit, i.e. $G^0$ est \uline{abélien} sur $S$, et
que la deuxième l'est sss $Q$ est plat sur $S$. / Cela prouve ainsi
que dans le cas $S=\mathrm{Spec}(V)$, on a (iv) $\Longrightarrow$ (i). C'est aussi
vrai dès que $S$ est seulement réduit, En effet, on sait cons%
% margin note written sideways
\margin{truire comme sous-schéma du schéma de Hilbert de $G/S$, un schéma $S'\to S$ résolvant le Pb universel correspondant à la recherche d'un $A$. Alors $S'\unsure{-|\to}S$ est un mono, et surjectif, prouvons que c'est un isom, i.e. \uline{propre}. Le critère valuatif nous ramène \unsure{au cas|encore à} $S=\mathrm{Spec}(V)$, où les remarques précédentes s'appliquent. D'où résulte l'existence du $A$, donc du $T=G/A$. Enfin $T$ est plat sur $S$ en vertu de (iv), [\ldots{} critère \unsure{val.|util.} \ldots].}
% archive sheet number: 195
